\documentclass[protokoll]{dvd}

\KOMAoptions{
    headwidth = 18cm,
    footwidth = 18cm,
}

% Styrelsen
\def\ordf{Daniell Cole}
\def\vice{\vak}
\def\samo{\vak}
\def\kassa{Owais Tabbaa}
\def\sekr{\vak}
\def\vak{\emph{vakant}}

% mötesinfo

\def\nr{2}
\def\yr{2026}
\def\ymd{2026-05-13}
\def\starttime{12:35}


\begin{document}

\title{Styrelsemöte \nr}
\subtitle{\yr}
\author{Styrelsen}
\date{\ymd}


\textbf{Datum:} \csname @date\endcsname\\
\textbf{Tid:} \starttime\\
\textbf{Plats:} Styrelserummet\\
\textbf{Styrelsemedlemmar:}
\begin{närvarande_förtroendevalda}
    \förtroendevald{Ordförande}{\ordf}{Ja}
    \förtroendevald{Vice Ordförande}{\vice}{---}
    \förtroendevald{Kassör}{\kassa}{Ja}
    \förtroendevald{SAMO}{\samo}{---}
    \förtroendevald{Sekreterare}{\sekr}{---}
\end{närvarande_förtroendevalda}

\textbf{Övriga medlemmar:}
\begin{närvarande_medlemmar}
\medlem{Ida Kjellerstedt}[]
\medlem{Alice Olson}[]
\end{närvarande_medlemmar}
\section{Öppnande av möte}

Mötet öppnades av \ordf\ kl \starttime

\section{Runda bordet}

Runda bordet innebär att varje person berättar hur de känner sig.
Man kan till exempel berätta att man är stressad på grund av en inlämning, irriterad på sin granne, eller bara väldigt glad därför att man ligger i fas med plugget.


\section{Formalia}

\subsection{Styrelsens beslutbarhet}

\blockquote[7 kap. 5 \S~första stycket i stadgan][]{%
    Styrelsen är endast beslutsmässig då samtliga styrelsemedlemmar har fått kallelsen till styrelsemötet och minst hälften av styrelsemedlemmarna är närvarande.
    Ordförande eller vice ordförande måste vara närvarande när beslut tas.
}

\blockquote[7 kap. 11 \S~första stycket i stadgan]{
Styrelsens sammanträden måste kallas minst två läsdagar innan sammanträdet. }

\blockquote[7 kap. 11 \S~sista stycket i stadgan]{
Kallelse till sammanträde får endast ske när protokollet från tidigare sammanträde är justerat och tillgängligt för medlemmarna.}

\subsubsection*{Förslag}

\begin{attsatser}
    \item Styrelsen har uppnått kraven i 7 kap. 5 § första stycket i stadgan och är därmed beslutbar.
\end{attsatser}
\subsubsection*{Beslut}
\begin{attsatser}
    \item förslaget till beslut bifalles
\end{attsatser}


\subsection{Fastställande av mötesschema}

För att styrelsen ska kunna fatta ett styrelsebeslut eller protokollföra en diskussion behöver punkten i mötesschemat där styrelsen ska fatta beslut vara inlagd eller föras in i mötesschemat senast vid den här punkten.

\subsubsection*{Förslag}

\begin{attsatser}
    \item 
    mötesschemat fastställs utan några förändringar.
    %punkt läggs till på mötesschema
\end{attsatser}
\subsubsection*{Beslut}
\begin{attsatser}
    \item förslaget till beslut bifalles
\end{attsatser}

\subsection{Val av mötessekreterare}

\subsubsection*{Förslag}
\begin{attsatser}
    \item \ordf\ väljs till mötessekreterare då det behövs en protokolljusterare men varken ordförande eller sekreterare kan får vara justerare.
\end{attsatser}
\subsubsection*{Beslut}
\begin{attsatser}
    \item förslaget till beslut bifalles
\end{attsatser}

\subsection{Val av protokolljusterare}

Protokolljusterare har till uppgift att kontrollera att protokollet i slutändan reflekterar de faktiska besluten och diskussionerna som fördes under mötet.
Utöver protokolljusteraren så ska mötesordförande och mötessekreteraren signera protokollet.
Vid styrelsemöten ska det endast vara en justerare.
Mötesordförande och mötessekreteraren kan inte vara justerare.

\subsubsection*{Förslag}
\begin{attsatser}
    \item \kassa\ väljs till protokolljusterare
\end{attsatser}
\subsubsection*{Beslut}
\begin{attsatser}
    \item förslaget till beslut bifalles
\end{attsatser}

\section{Rapport} 
\subsection{Ordförande}
Har uppdaterat stadgan utefter beslut som röstades igenom på stämman förra året.

Gått på flertal möten med göta. 

Att vi ska bli organisationens fullmakt är på väg, 

Fotograferingen har jag pratat med Sarah om, Ester kommer skicka en faktura. 


\subsection{Kassör}
Väntar på att fullmakt ska gå igenom. 


\newpage


\section{Diskussionspunkter}
\subsection{Frimärken för att posta till banken}
Styrelsen behöver 8 st frimärken, 22 kr styck, för att kunna posta papper till banken.  För detta äskar vi 
Styrelsen har beslutat att bevilja 176 kr till att köpa frimärken. 

\section{Besultspunkter}
\subsection{Äskan: Vattenkokare}
Det har tillkommit efterfrågan på ny vattenkokare till Monadens kök. Den som finns börjar dra på åren och ser ut som den är bränd i botten. Det är på tiden att den byts ut.

Föreslagen modell: Wilfa WKD-2200S, kostar 500kr.

\subsubsection*{Förslag till beslut}
\begin{attsatser}
    \item äskan för den föreslagna modellen på vattenkokare godkänns och skall köpas in när överföring av fullmakt färdigställts. 
\end{attsatser}
\subsubsection*{Beslut} 
\begin{attsatser}     
\item attsats 1 bifalles. 
\end{attsatser}


\subsection{Föreningsfotografering}
Esther är en professionell fotograf som har erbjudit att ta våra bilder för 4 000 kr (3 000  kr + 1500 kr reseersättning). 
Detta är ett otroligt bra pris för en professionell fotograf med den mängd och typ av bilder i åtanke. 

\subsubsection*{Förslag till beslut}
\begin{attsatser}
    \item detta erbjudande godtas.
\end{attsatser}
\subsubsection*{Beslut}
\begin{attsatser}
    \item attsats 1 bifalles.
\end{attsatser}

\subsection{Betalning från göta (kårkällar'n) har gått till fel konto}
Göta hade under flera år haft fel konto inlagt i deras system. Detta har bara påverkat när kårkällar'n skulle betala ut till oss direkt då äskningar betalas direkt till personen som gjort utlägget. Albin har tillgång till det konto där pengarna hamnat. Jag har redan haft samtal med honom om detta och vi behöver skicka en officiell förfrågan till Albin att skicka över pengarna. 
\subsubsection*{Förslag till beslut}
\begin{attsatser}
    \item ordförande ber Albin Otterhäll att överför pengarna i fråga.
\end{attsatser}
\subsubsection*{Beslut} 
\begin{attsatser}     
\item attsats 1 bifalles. 
\end{attsatser}



\newpage
\section{Avslutande av möte}

\subsection{Nästa möte}
2026-05-27 - 12:12

\subsection{Mötets avslutande}
Mötet avslutades kl 13:13 

\styrelsesignaturer

\end{document}